%% P3 — Vietnam
%% Target journal: Journal of Asia Business Studies (JABS), Emerald Publishing
%% Article class: standard article (Emerald accepts article class submissions)
%% Compile: pdflatex → biber → pdflatex × 2  |  or: latexmk -pdf p3_jabs_submission.tex

\documentclass[12pt,a4paper]{article}

%% ---- Packages ----
\usepackage{amsmath,amssymb,amsthm}   % math environments
\usepackage{graphicx}
\usepackage{booktabs}                  % professional tables (\toprule, \midrule, \bottomrule)
\usepackage{setspace}
\doublespacing
\usepackage{geometry}
\geometry{margin=2.5cm}
\usepackage[style=apa,backend=biber]{biblatex}
\addbibresource{../../references/p3_references.bib}
\usepackage{hyperref}
\hypersetup{colorlinks=true, linkcolor=black, citecolor=black, urlcolor=blue}
\usepackage{lmodern}
\usepackage{microtype}

%% ---- Shared macros ----
%% Shared macros for PhD dissertation papers
%% Internationalization-Performance relationship studies (P3, P4, P5, P6)

%% --- Key variables (italic, math mode) ---
\newcommand{\lnLP}{\ln(\mathrm{LP})}
\newcommand{\FSTS}{FSTS}
\newcommand{\FSTSc}{FSTS_{c}}
\newcommand{\FSTScsq}{FSTS_{c}^{2}}
\newcommand{\FSTSccb}{FSTS_{c}^{3}}
\newcommand{\TCI}{\mathrm{TCI}}
\newcommand{\TCIz}{\mathrm{TCI}_{z}}
\newcommand{\DAI}{\mathrm{DAI}}
\newcommand{\DAIz}{\mathrm{DAI}_{z}}
\newcommand{\IMR}{\mathrm{IMR}}

%% --- Model notation ---
\newcommand{\bvec}[1]{\boldsymbol{#1}}
\newcommand{\bX}{\bvec{X}}
\newcommand{\bgamma}{\bvec{\gamma}}
\newcommand{\bcontrols}{\bgamma'\bX_{i}}

%% --- ICRV regimes ---
\newcommand{\ICRVadv}{\mathrm{ICRV}_{\mathrm{Adv}}}
\newcommand{\ICRVem}{\mathrm{ICRV}_{\mathrm{Em}}}
\newcommand{\ICRVfr}{\mathrm{ICRV}_{\mathrm{Fr}}}

%% --- Fixed effects ---
\newcommand{\alphaj}{\alpha_{j}}   % country FE
\newcommand{\deltat}{\delta_{t}}   % year FE

%% --- Turning point shorthand ---
\newcommand{\TP}{\mathrm{TP}}

%% --- Statistical operators ---
\newcommand{\pval}{p}
\newcommand{\betacoef}{\hat{\beta}}


%% ---- Title / Author (blinded for review) ----
\title{Revisiting the Internationalisation--Performance Relationship in an
  Emerging Market: The Roles of Technological Capability and Foundational
  Digital Adoption}
\author{[Blinded for review]}
\date{}

% ============================================================
\begin{document}
\maketitle

%% ---- ABSTRACT ----
\begin{abstract}
\noindent\textbf{Purpose} --- This study revisits the
internationalisation--performance (\(I \to P\)) relationship in Vietnam across
three survey waves, testing whether technological capability moderates
export-intensity effects and how a Tier-1 digital indicator evolves across
institutional generations.

\medskip
\noindent\textbf{Design/Methodology/Approach} --- Three waves of World Bank
Enterprise Survey data (Vietnam 2009, 2015, 2023; \(N = 2{,}958\)) are used to
estimate wave-specific and pooled OLS models with HC1 robust standard errors,
quadratic export-intensity terms, and capability interactions.
\(\mathrm{TCI}\) (quality certification, foreign-licensed technology) is the
primary construct; a single-item website indicator (Tier-1 only) is retained
as a digital-presence control.

\medskip
\noindent\textbf{Findings} --- An inverted-\(U\) between export intensity and
labour productivity holds across all three waves, with turning points clustered
between 39\% and 46\% --- a structurally durable 7-percentage-point band
sustained across 14 years of institutional change.  The Lind--Mehlum test
rejects monotonicity in all waves (\(p = .006, .009, .013\) and pooled
\(p < .001\)).  \(\TCIz\) is positively associated with productivity in all
waves (pooled \(\hat{\beta} = 0.179\), \(p < .001\)) and moderates the \(I
\to P\) curvature.

\medskip
\noindent\textbf{Originality/Value} --- The inverted-\(U\) nonlinearity
reflects a participation-margin effect at the export-participation barrier
rather than a continuous intensity-saturation curve.  The threshold cluster
(39--46\% \(\FSTS\) band) is institutionally embedded --- shifting the debate
from ``does internationalisation help?'' to ``under what institutional
conditions does the threshold shift?''

\medskip
\noindent\textbf{Keywords:} internationalisation--performance; emerging
markets; technological capability; threshold durability; Vietnam; firm
productivity.

\medskip
\noindent\textbf{JEL:} F23; O33; D22; L25; O53.
\end{abstract}

\newpage

% ============================================================
\section{Introduction}
\label{sec:intro}

Vietnam offers an analytically valuable setting for revisiting the
internationalisation--performance (\(I \to P\)) relationship because firms
expand abroad under conditions of institutional transition, uneven capability
accumulation, and rapidly changing digital infrastructure.  A long tradition of
research supports the view that nonlinearity is a central feature of this
relationship \parencite{Lu2004, Hennart2007, Marano2016}.

Three institutional turning points shape the observation window: WTO accession
(2007), the mid-transition CPTPP/EVFTA period (2015), and the National Digital
Transformation Programme post-2020 (2023 wave).

% ============================================================
\section{Theory and Hypotheses}
\label{sec:theory}

\subsection{Inverted-\textit{U} Nonlinearity (H1)}

Building on Uppsala model learning logic \parencite{Johanson1977,Johanson2009}
and organisational complexity arguments \parencite{Lu2004}, we hypothesise:

\begin{quote}
\textbf{H1:} The relationship between export intensity (\(\FSTS\)) and labour
productivity (\(\lnLP\)) follows an inverted-\(U\) shape, with a maximum in the
intermediate range.
\end{quote}

\subsection{TCI Moderation (H2)}

Resource-based view \parencite{Barney1991} predicts that firms with higher
technological capability sustain the performance premium at higher
internationalisation levels:

\begin{quote}
\textbf{H2:} Higher \(\TCI\) raises the productivity level (positive level
shift) and moderates the curvature of the \(I \to P\) relationship.
\end{quote}

\subsection{DAI Proxy-Obsolescence (H3/H4)}

The Tier-1 digital-presence indicator (website ownership) follows a
proxy-obsolescence trajectory as adoption diffuses to near-universal levels:

\begin{quote}
\textbf{H3/H4:} The \(\DAI\) (Tier-1 website) coefficient evolves from
positive (2009) to null (2015) to negatively interactive (2023), consistent with
Tier-1 proxy obsolescence as website adoption diffuses toward universal levels ---
not evidence of digital adoption moderation in the capability sense.
\end{quote}

% ============================================================
\section{Methodology}
\label{sec:methods}

\subsection{Data}

World Bank Enterprise Surveys (WBES): Vietnam 2009 (full frame \(n = 1{,}074\)),
2015 (\(n = 1{,}001\)), 2023 (\(n = 883\)).  Analytic sample after listwise deletion
on focal variables: 2009 \(n = 989\), 2015 \(n = 956\), 2023 \(n = 1{,}013\);
pooled \(N = 2{,}958\).  All monetary values converted to 2017 constant USD PPP.

\subsection{Variables}

\paragraph{Dependent variable.}
Labour productivity: \(\lnLP = \ln(\text{annual revenue PPP} / \text{permanent
workers})\).

\paragraph{Internationalisation.}
Foreign Sales to Total Sales, mean-centred: \(\FSTSc = \FSTS - \bar{F}\).

\paragraph{Technological Capability Index.}
Within-wave standardised composite of quality certification (\texttt{b8}) and
foreign-licensed technology (\texttt{e6}):
\begin{equation}
  \TCIz = \frac{1}{2}\left[z(\text{quality\_cert}) + z(\text{for\_lic\_tech})\right].
  \label{eq:tci}
\end{equation}

\paragraph{Digital Adoption Index.}
Single-item Tier-1 binary: website presence (\texttt{c22b}), standardised.

\subsection{Model Specification}

\paragraph{M0 --- Linear baseline.}
\begin{equation}
  \lnLP_{i} = \beta_{0} + \beta_{1}\,\FSTS_{i}
    + \bcontrols + \alphaj + \deltat + \varepsilon_{i}
  \label{eq:M0}
\end{equation}

\paragraph{M1 --- Quadratic inverted-\textit{U} (H1).}
\begin{equation}
  \lnLP_{i} = \beta_{0} + \beta_{1}\,\FSTSc_{i} + \beta_{2}\,\FSTScsq_{i}
    + \bcontrols + \alphaj + \deltat + \varepsilon_{i}
  \label{eq:M1}
\end{equation}
Turning point: \(\TP = -\beta_{1}/(2\beta_{2})\).  Confirmed by
Lind--Mehlum \textit{utest} \parencite{LindMehlum2010}.

\paragraph{M2 --- Cubic S-curve (H1 extended).}
\begin{equation}
  \lnLP_{i} = \beta_{0} + \beta_{1}\,\FSTSc_{i} + \beta_{2}\,\FSTScsq_{i}
    + \beta_{3}\,\FSTSccb_{i}
    + \bcontrols + \alphaj + \deltat + \varepsilon_{i}
  \label{eq:M2}
\end{equation}

\paragraph{M3 --- TCI moderation (H2).}
\begin{align}
  \lnLP_{i} &= \beta_{0} + \beta_{1}\,\FSTSc_{i} + \beta_{2}\,\FSTScsq_{i}
    \notag \\
    &\quad + \beta_{3}\,\TCIz_{i}
    + \beta_{4}\,(\FSTSc_{i} \times \TCIz_{i})
    + \beta_{5}\,(\FSTScsq_{i} \times \TCIz_{i})
    \notag \\
    &\quad + \bcontrols + \alphaj + \deltat + \varepsilon_{i}
  \label{eq:M3}
\end{align}
Joint moderation test: \(F\)-test \(H_{0}\colon \beta_{4} = \beta_{5} = 0\).

\paragraph{M4 --- DAI moderation (H3/H4).}
\begin{align}
  \lnLP_{i} &= \beta_{0} + \beta_{1}\,\FSTSc_{i} + \beta_{2}\,\FSTScsq_{i}
    \notag \\
    &\quad + \beta_{3}\,\DAIz_{i}
    + \beta_{4}\,(\FSTSc_{i} \times \DAIz_{i})
    + \beta_{5}\,(\FSTScsq_{i} \times \DAIz_{i})
    \notag \\
    &\quad + \bcontrols + \alphaj + \deltat + \varepsilon_{i}
  \label{eq:M4}
\end{align}

\paragraph{Controls.}
\(\bX_{i}\) includes: \(\ln(\text{employees})\), firm age,
foreign-ownership dummy (\(\geq 10\%\)), sector dummies.
HC1 robust standard errors throughout \parencite{LongErvin2000}.

% ============================================================
\section{Results}
\label{sec:results}

\subsection{Descriptive Statistics}

[Table~1 --- Descriptives and correlations (mean, SD, min, max,
Pearson~\(r\) by wave)]

\subsection{Baseline and Nonlinearity Tests}

The quadratic term is negative and significant in all three waves
(pooled \(\hat{\beta}_{2} = -0.441\), \(p < .001\)).  Turning points:
2009 = 46.2\%, 2015 = 39.3\%, 2023 = 41.6\%, pooled = 39.7\%.
The Lind--Mehlum \textit{utest} confirms an interior maximum in all panels.

\subsection{TCI Moderation}

\(\TCIz\) is positively associated with \(\lnLP\) in all waves
(\(\hat{\beta}_{3} = 0.215, 0.128, 0.123\)) and pooled
(\(\hat{\beta}_{3} = 0.179\), \(p < .001\)).
Joint moderation \(F\)-tests: \(p = .040, p > .100, p = .027, p = .003\)
for waves 2009, 2015, 2023, and pooled, respectively.

\subsection{DAI Proxy Obsolescence}

Baseline DAI coefficients: 2009 \(\hat{\beta} = 0.175\) (positive),
2015 \(\hat{\beta} = -0.044\) (null), 2023 \(\hat{\beta} = 0.095\) (positive).
The \(\FSTSc \times \DAIz\) interaction is significant only in 2023
(\(\hat{\beta}_{4} = -0.912\), \(p = .043\)), consistent with proxy
obsolescence as website adoption diffuses to 49.8\% of firms.

\subsection{Robustness}

[Table~3 --- Alternative estimators (random effects, GMM), alternative
turning-point methods, exporter-only sub-sample, winsorisation sensitivity]

The quadratic term on the exporter-only sub-sample loses significance
(\(\hat{\beta} = -0.200\), \(p = .660\)), confirming that the inverted-\(U\)
is driven by the participation margin rather than within-exporter intensity
saturation.

% ============================================================
\section{Discussion}
\label{sec:discussion}

The structural durability of the 39--46\% turning-point band across 14 years
suggests the threshold is institutionally embedded rather than firm-specific.
This reframes the I--P debate: the relevant constraint in a transitional
economy binds at the export-participation barrier, not within the intensive
margin.

The proxy-obsolescence trajectory of the Tier-1 digital indicator provides
methodological guidance for future research: single-item digital measures
require periodic re-validation as adoption levels change.

% ============================================================
\section{Conclusion}
\label{sec:conclusion}

Across three survey waves in Vietnam, the inverted-\(U\) between export
intensity and labour productivity is structurally durable (39--46\% band),
primarily reflects a participation-margin step function, and is positively
moderated by technological capability.  The website indicator follows a
proxy-obsolescence arc rather than moderation by digital adoption capacity.

\paragraph{Limitations and future research.}
The DAI measure is limited to Tier-1 (website presence), a measurement
constraint imposed by WBES Vietnam questionnaire design.  Future research
should incorporate Tier-2 transaction-enabling digital indicators
(e-payment, \texttt{k33}/\texttt{k38}) when available across waves.

% ============================================================
\printbibliography

\end{document}
